\section{Conclusion}
\label{sec:conclusion}

We rebuilt a multi-source six-class chest radiograph benchmark from official sources with labelling rules fixed before training, patient- and duplicate-grouped splits (separation proven only where patient identity is known) and a held-out external cohort, and used it to examine source shortcuts in a standard DenseNet-121. Internal test accuracy was \RTestAcc\% (Wilson 95\% interval \RTestAccLo--\RTestAccHi\%), yet only \RExtSpecK{} of \RExtSpecN{} normal external radiographs were called Normal and \RExtTBK{} of \RExtTBN{} TB radiographs were named Tuberculosis, consistently across three seeds.

In frozen features, source information was recoverable from every spatial-frequency band, from \RSpecMainTestPassASrcBA\% to \RSpecMainTestPassFSrcBA\% balanced accuracy against a chance level of 16.7\%, and no single band was necessary for it. To our knowledge, within the searches described, this band-by-band localisation of source versus disease information has not been reported for chest radiographs; its outcome is that the signal is not localised. Linear erasure lowered the recoverable source information from \RInlpImgTestSrcBAKZero\% to \RInlpImgTestSrcBAKStar\%, not to chance, and lowered disease scores too, consistent with entanglement. A linear head trained without one source scored \RLosoImgCohenGap{} to \RLosoImgTBXkGap{} points below its in-distribution reference on four of five held-out sources. The planned band-stop retraining was not run because the rule fixed in advance returned \RBsDecision, an outcome we record as part of the design. These results come from linear methods on frozen features of one benchmark and one run, carry the patient-identity caveat, and are indirect evidence about what a trained model uses.

For multi-source chest radiograph work the practical lesson is modest: document provenance, fix label rules in advance, group splits by patient and content, report per source, evaluate on an unseen hospital with sensitivity and specificity together, and note that these results give no reason to expect frequency filtering or linear erasure of frozen features to replace source-balanced data, which we did not test. We make no claim of clinical validity.
